\chapter{Флоке-устойчивость минимального двухузлового компакттона}
\label{ch:v6-spectral-transition-discrete-compacton-stability-quantization-gate}

\section{Строгая постановка после гейта существования}

Предыдущая глава построила точную четырёхпериодическую орбиту
$\Psi_\star$ при $\kappa=2\pi$ и доказала отсутствие утечки на самой
орбите. Это ещё не равносильно устойчивости: нужно проверить производную
возвратной карты
\begin{equation}
 \mathcal M=D_{\mathbb R}F_{2\pi}^{,4}(\Psi_\star).
 \label{eq:v6-compacton-real-monodromy}
\end{equation}
Производная берётся над вещественными и мнимыми частями отдельно, поскольку
составной дублет $H_{\rm eff}=\ell\overline e$ содержит комплексное
сопряжение и карта не является голоморфной.

Этот тест продолжает строгий ledger Тома V: глава
\path{version5_transition_primitive_scientific_language_gate} требовала
отличать устойчивый дефект композиции от одной точной орбиты, а глава
\path{version6_spectral_transition_discrete_compacton_existence_gate}
оставила именно монодромию открытым условием. Первичная литература по
локальным унитарным автоматам и нелинейным квантовым блужданиям
подтверждает допустимость такого языка, но не фиксирует ни $\kappa$, ни
устойчивость данного проектного профиля
(\path{arXiv:hep-th/9304070}, \path{arXiv:1902.02017},
\path{arXiv:2308.01014}).

\section{Полный вещественный спектр}

Центральной разностью с шагом $2\cdot10^{-7}$ вычислена полная матрица
\eqref{eq:v6-compacton-real-monodromy} на периодических решётках трёх
размеров. Получено:
\begin{center}
\begin{tabular}{@{}rrrrr@{}}
\toprule
$N$ & вещественный размер & $\rho(\mathcal M)$ & $|\lambda|>1+10^{-5}$ & $\|\mathcal M\|_2$ \\
\midrule
8  & 96  & $1.000000001006$ & 0 & $6.438500961$ \\
12 & 144 & $1.000000000931$ & 0 & $6.438500961$ \\
16 & 192 & $1.000000000891$ & 0 & $6.438500961$ \\
\bottomrule
\end{tabular}
\end{center}
Остаток четырёхшагового возврата равен $2.61\cdot10^{-15}$. Рост
спектрального радиуса не усиливается с объёмом и лежит на уровне ошибки
численного дифференцирования. Расширяющей флоке-моды не найдено.

\section{Удаление симметрий}

Проверены касательные к общей фазе, трём слабым вращениям Паули и фазе
правого хирального конца. Из пяти записанных направлений независимы четыре;
максимальный остаток условия $\mathcal Mv=v$ равен
$1.13\cdot10^{-9}$. После ортогонального удаления этого четырёхмерного
пространства остаётся вещественная матрица размера $188$ с
\begin{equation}
 \rho(\mathcal M_{\rm red})=1.000000000901,
 \qquad
 \#\{|\lambda|>1+10^{-5}\}=0.
 \label{eq:v6-compacton-reduced-radius}
\end{equation}
Следовательно, положительный результат не создаётся одной фазовой или
калибровочной нулевой модой.

Однако в редуцированном спектре остаются $128$ мультипликаторов на
единичной окружности. Это ожидаемый нейтральный сектор безмассового
вакуумного сдвига. Компактон спектрально устойчив, но не является
асимптотическим аттрактором: малое излучение не обязано возвращаться в
ядро.

\section{Ненормальность и конечные возмущения}

Матрица монодромии ненормальна. Максимальные сингулярные числа её степеней
для $n=1,2,4,8,16,32,64,128$ равны соответственно
\begin{equation}
 6.4385,\ 8.9969,\ 11.8392,\ 4.6576,\
 6.4385,\ 8.9969,\ 11.8392,\ 4.6576.
\end{equation}
Это допускает заметное переходное усиление, но в проверенном окне оно
ограничено и повторяется, а не растёт экспоненциально.

Дополнительно взято одно фиксированное случайное комплексное направление
единичной нормы на решётке из $256$ узлов. После добавления возмущения
нормы $\delta$, общей нормировки и $80$ шагов вероятность в исходном
двухузловом ядре равна:
\begin{center}
\begin{tabular}{@{}rr@{}}
\toprule
$\delta$ & $P_{\rm core}(80)$ \\
\midrule
$10^{-4}$ & $0.9999999909$ \\
$10^{-3}$ & $0.9999990896$ \\
$10^{-2}$ & $0.9999073772$ \\
$0.03$ & $0.9990061304$ \\
$0.05$ & $0.9954525401$ \\
$0.06$ & $0.9879834021$ \\
$0.07$ & $0.5952017862$ \\
$0.08$ & $0.4218826447$ \\
$0.09$ & $0.0119538725$ \\
$0.10$ & $0.0091284718$ \\
\bottomrule
\end{tabular}
\end{center}
Для малых $\delta$ излучённая вероятность имеет порядок $\delta^2$.
Для данного воспроизводимого направления резкий уход начинается между
$0.06$ и $0.07$. Этот интервал является свойством протокола, а не
универсальной константой: он доказывает наличие малого устойчивого окна и
одновременно опровергает глобальную устойчивость.

\begin{theorem}[Численная флоке-устойчивость ветви $2\pi$]
На периодических решётках $N=8,12,16$ полная и симметрийно редуцированная
вещественные монодромии минимального компакттона не имеют
мультипликаторов с модулем выше $1+10^{-5}$. Малые контролируемые
возмущения сохраняют почти всю норму в двухузловом ядре на протяжении
$80$ шагов. Поэтому квантованная ветвь $\kappa=2\pi$ проходит первый
локальный тест линейной устойчивости.
\end{theorem}

\begin{nogocriterion}
Результат нельзя называть доказательством полной устойчивости частицы.
Сохраняются обширный нейтральный радиационный сектор, ненормальное
переходное усиление, конечнопороговый уход в выбранном протоколе и
отсутствие аналитической нелинейной теоремы на бесконечной решётке.
Физический шаг решётки, масса и наблюдаемая карта также не выведены.
\end{nogocriterion}

\section{Следующий гейт}

Следующий гейт
\path{version6_spectral_transition_discrete_compacton_physical_scale_map_gate}
должен проверить, выводятся ли физическая длина двухузлового ядра и
квазиэнергия фазы $\pm i$ из уже существующих масштабных соотношений S2T,
не вводя новый внешний шаг решётки.

Нулевая гипотеза: безразмерный автомат оставляет шаг $a$ свободным, поэтому
размер $2a$ и энергия $\pi/(2\Delta t)$ не являются слепыми
предсказаниями. Стоп-критерий: если ни один существующий масштаб проекта
не фиксирует $a$ и $\Delta t$ одним родительским законом, $2\pi$ остаётся
квантованием существования, но не физической массой.

\begin{protocol}
\begin{enumerate}
 \item точность четырёхшаговой орбиты: \textbf{$2.61\cdot10^{-15}$};
 \item расширяющие полные флоке-моды: \textbf{не найдены};
 \item расширяющие редуцированные моды: \textbf{не найдены};
 \item независимые симметрийные касательные: \textbf{четыре};
 \item нейтральный радиационный сектор: \textbf{сохраняется};
 \item максимальное проверенное переходное усиление:
       \textbf{$11.83923$};
 \item малое нелинейное окно: \textbf{найдено численно};
 \item глобальная нелинейная устойчивость: \textbf{не доказана};
 \item ветвь $\kappa=2\pi$: \textbf{проходит локальный тест устойчивости};
 \item R4: \textbf{частично закрыт};
 \item R5 и физический масштаб: \textbf{открыты}.
\end{enumerate}
\end{protocol}

Машинный сертификат сохранён в
\path{s2t_v6_spectral_transition_discrete_compacton_stability_quantization_gate_results.json}.